\documentclass[11pt,a4paper]{article}
\usepackage[T1]{fontenc}
\usepackage{lmodern,microtype,amsmath,amssymb,amsthm}
\usepackage[margin=25mm]{geometry}
\usepackage[colorlinks=true,urlcolor=blue,linkcolor=blue]{hyperref}
\setlength{\emergencystretch}{3em}
\newtheorem{theorem}{Theorem}
\newtheorem{corollary}[theorem]{Corollary}
\newcommand{\PP}{\mathbb P}
\title{Zero branches preserve the entire class of causal seed laws}
\author{Research note for the finite-Shannon-theorem project}
\date{11 October 2026}
\begin{document}
\maketitle
\begin{abstract}
For a locally finite controlled tree with nonnegative normalized kernels,
we prove that positive-support reduction preserves exactly which probability
spaces can serve as causal seeds, simultaneously at every finite horizon.
The seed is unchanged; only the decoder is transformed. The analogous
statement holds for ordinary policy-indexed sources. Consequently the
positive-kernel infinite-horizon theorems transfer to the original tree,
with its zero-probability branches retained and its decoder total there.
This is a mathematical extension, not a new Lean certificate or a claim
of historical priority. The underlying support-reduction idea is standard.
\end{abstract}

\section{Model and an exact equivalence}
At each finite history $h$, let $A(h)$ be a finite nonempty action set;
for $a\in A(h)$ let $Y(h,a)$ be a finite nonempty output set, with
\[
 p(y\mid h,a)\ge0,\qquad \sum_{y\in Y(h,a)}p(y\mid h,a)=1.
\]
Neither stationarity nor a uniform positive lower bound is assumed.
A terminal history, if present, can be continued by a deterministic dummy
action and output. The set of finite histories is countable.

Call a history \emph{positive} if all transition probabilities along it
are positive, and let $H^+$ be their set, including the empty history.
The auxiliary tree $T^+$ keeps the same actions at histories in $H^+$
and replaces $Y(h,a)$ by
$Y^+(h,a)=\{y:p(y\mid h,a)>0\}$. Its transition probabilities are
unchanged: there is no renormalization. Write $j$ for inclusion of its
histories into those of the original tree $T$.

On a probability space $(\Omega,\mathcal F,\mu)$, a causal decoder is a
measurable response $D(\omega,h,a)\in Y(h,a)$ for every history and action.
It is exact if every deterministic adaptive policy has the prescribed
finite transcript laws. Measurability means measurability of each finite
output coordinate. An ordinary source instead has separate measurable
transcript readouts for each finite horizon and policy, with the same
exact marginals and without consistency requirements between readouts.

\begin{theorem}[Preservation of the seed space]\label{thm:seed}
For every probability space $(\Omega,\mathcal F,\mu)$,
\[
 \boxed{\quad
 (\Omega,\mu)\text{ admits an exact causal decoder for }T
 \iff
 (\Omega,\mu)\text{ admits one for }T^+ .\quad}
\]
The equivalence holds both at any fixed horizon and simultaneously at
all finite horizons. It also holds with ``causal decoder'' replaced by
``ordinary source.'' All conversions leave $\Omega$, $\mathcal F$, and
$\mu$ unchanged. On total strategy spaces the extension map is an
injective measurable section of the restriction map.
\end{theorem}

\begin{proof}
Choose a positive default output $y_*(h,a)\in Y^+(h,a)$ at every original
history and action; normalization ensures existence. Put
\[
 r_{h,a}(y)=
 \begin{cases}y,&p(y\mid h,a)>0,\\y_*(h,a),&p(y\mid h,a)=0.\end{cases}
\]
For a full decoder $D$ define its reduced decoder by
\[
 RD(\omega,h,a)=r_{j(h),a}(D(\omega,j(h),a)).
\]
Conversely, for a reduced decoder $D^+$, define a \emph{total} full decoder
\[
 ED^+(\omega,h,a)=
 \begin{cases}
 D^+(\omega,j^{-1}(h),a),&h\in H^+,\\
 y_*(h,a),&h\notin H^+.
 \end{cases}
\]
Both maps are measurable coordinatewise, and $RED^+=D^+$ pointwise.
In particular $E$ is injective. The definitions do not depend on a horizon.

Suppose first that $D^+$ is exact. For any full policy $\pi$, restrict
its choices to positive histories, obtaining $\pi^+$. Running $ED^+$
under $\pi$ stays in $H^+$ for \emph{every} seed value and is precisely
the included run of $D^+$ under $\pi^+$. Under the prescribed full law,
any path leaving $H^+$ has probability zero, while every positive path
has its unchanged product probability. Thus $ED^+$ is exact.

Now suppose $D$ is exact. Extend a reduced policy $\pi^+$ to a full
policy $\pi$ by arbitrary legal choices outside $H^+$. For each finite
horizon, the $D$-run under $\pi$ has probability zero of taking a zero
transition, by its exact transcript law. Until such a transition, the
$D$-run and the $RD$-run coincide. Therefore the reduced run has exactly
the required law. This proves exactness at every horizon using the
single decoder $RD$ above; it is not an argument that chooses separate
seeds or decoders for each horizon.

There is also a common full-measure set on which $D$ never chooses a zero
output at any history reachable under any policy. For each finite path
ending in a zero transition, consider the event that $D$ agrees with
all its specified outputs. A policy that follows the specified actions
on that path shows that this event has measure zero. There are countably
many such paths, so their union is null. No uncountable union over
infinite policies is needed. This statement concerns reachable histories,
not arbitrary unused coordinates of a strategy.

For ordinary sources, reduced-to-full conversion restricts a full policy,
uses its reduced readout, and includes the resulting transcript.
In the other direction, extend the reduced policy and use the full
readout. A positive transcript is restricted; replace a nonpositive one
by any fixed positive transcript compatible with the reduced policy.
The replacement event has measure zero, and the readout remains measurable.
This construction works separately at each readout index, exactly as the
definition of an ordinary source permits. Throughout, the seed law is
literally unchanged.
\end{proof}

\section{Spectra and the infinite consequences}
For a fixed policy and horizon, support reduction preserves the multiset
of positive transcript masses. Every full policy restricts and every
reduced policy extends. Consequently, for every $n,t,K$,
\begin{equation}\label{eq:spectrum}
 F_{T_n}^*(t)=F_{T_n^+}^*(t),\qquad
 \inf_\pi A_K(P^T_{n,\pi})=\inf_\pi A_K(P^{T^+}_{n,\pi}).
\end{equation}
Here $A_K$ is the sum of the $K$ largest masses, padded with zeros if
necessary. Information distributions sum over positive transcript masses.
Zero-mass leaves contribute no probability at $+\infty$.

\begin{corollary}\label{cor:transfer}
The positive-tree infinite-horizon results extend to $T$ with all zero
branches retained. In particular:
\begin{enumerate}
\item For countable-seed Shannon cost, in a single consistent unit,
\[
 C_\infty(T)=\sup_n C_n(T).
\]
If finite, the optimum is attained by a countable seed on the original tree.
\item In bits, with $L_\infty=\sup_n L(T_n)$ and
$c_*=(1+\log_2e)/2$,
\[
 L_\infty\le C_\infty(T)\le L_\infty+c_*.
\]
An exact finite-entropy seed exists if and only if $L_\infty<\infty$.
\item For sources with measurable singletons, define
$q=\sup_K\inf_{n,\pi}A_K(P^T_{n,\pi})$ and
$F_\infty(t)=\inf_nF_{T_n}^*(t)$. The maximal ordinary and causal atomic
capacities equal $q$, the causal capacity is attained, and
\[
 q=\lim_{t\to\infty}F_\infty(t)=\sup_t F_\infty(t).
\]
An exact countably supported causal seed exists if and only if $q=1$.
\item If a countable ordinary exact source exists, there is one countable
causal seed $S$, exact at all horizons, for which simultaneously for all
$\alpha>0$,
\[
 R_\alpha(Z)\le O_\alpha(T)\le C_\alpha(T)
 \le H_\alpha(S)\le R_\alpha(Z)+d_\alpha,
\]
with the same constants as for the positive tree. Its min-entropy attains
both the ordinary and causal optima. Here $Z$ is the envelope-limit law,
including its possible mass $1-q$ at $+\infty$.
\end{enumerate}
\end{corollary}
\begin{proof}
Theorem~\ref{thm:seed} identifies, at every finite horizon and at the
infinite horizon, the sets of admissible countable seed probability laws.
Thus Shannon costs and all ordinary and causal R\'enyi costs agree with
those of $T^+$. Entropy attainment and the simultaneous choice of a seed
transfer by $E$, without altering that seed's weights.
Equation~\eqref{eq:spectrum} gives equality of $L_\infty$, $F_\infty$,
$q$, and the envelope-limit law, so all the displayed spectral bounds
follow from the positive-tree theorems.

For atomic capacity, apply the ordinary upper bound to the unchanged
source space using $R$. To transfer attainment as a measure on total
strategies, push the positive-tree attaining measure through $E$.
Because $E$ is injective, every target singleton has as preimage either
one singleton or the empty set; total singleton mass is preserved.
The pushed-forward law is exact by Theorem~\ref{thm:seed}. This also
justifies attainment in the strategy-space formulation, rather than
only on an abstract seed space.
\end{proof}

\section{Scope, dependencies, and what remains open}
The extension retains a total decoder on the original tree. The positive
tree is an auxiliary proof device, not a replacement of the final model.
No limiting perturbation $p\mapsto p+\varepsilon$ is used, so there is no
interchange of a positivity limit with an infinite horizon.

The input results are Theorems 5--8 of the repository README and its
atomic spectral identity; the new mathematical obligation is entirely
Theorem~\ref{thm:seed} and \eqref{eq:spectrum}. Relevant Lean modules are
\texttt{EntropyLimit}, \texttt{InfiniteEnvelope}, \texttt{AtomicCapacity},
\texttt{AtomicSpectrum}, and \texttt{RenyiLimitMeasure}. The repository
already has finite stationary support reduction in \texttt{ZeroKernel}
and \texttt{ZeroCost}; this note supplies the general horizon-independent
argument for history-dependent trees. It has not been formalized or
compiled in Lean.

A stronger claim that adjoining arbitrary strictly lower-entropy actions
changes every $C_n$ by at most a constant is \emph{not proved here}.
The bounded-excess theorem for a unique maximizing action does not alone
establish that claim when several actions maximize entropy.
The exact $G_*$, endpoint minimizers, optimal overhead constants, and the
equal-variance next-order problem remain unresolved by these arguments.

\end{document}
